\section{Cylinder geometry and causal realization}\label{sec:v5-realization}
Throughout this section $q$ and the data in
\eqref{eq:v5-branches}--\eqref{eq:v5-gibbs} are fixed. We suppress the
subscript $q$ in $V$ and $a$. A comparison constant never depends on a
word or a state budget.

\begin{lemma}[Distortion and orbit separation]\label{lem:v5-cylinder}
There is $D\ge1$ such that
\begin{equation}\label{eq:v5-distortion}
 \sup h_w'\le D\inf h_w',\qquad
 D^{-1}r_u r_v\le r_{uv}\le r_u r_v.
\end{equation}
If $x,y\in K$ have longest common symbolic prefix $w$, then
\begin{equation}\label{eq:v5-orbitseparation}
 c_1V(w)\le\|Z_q(x)-Z_q(y)\|^2\le V(w),
 \qquad c_1=(1-q)g^2D^{-2}.
\end{equation}
For arbitrary words $u,v$,
\begin{align}
 V(uv)&\le r_v^2\bigl(V(u)-q^{|u|}\bigr)+q^{|u|}V(v),
       \label{eq:v5-concat}\\
 V(uv)&\le V(u)V(v),\qquad a(uv)\le a(v).
       \label{eq:v5-suffixdomination}
\end{align}
There exists $0<a_0<\rho_q$ such that each one-symbol child
$wi$, $1\le i\le b$, satisfies
\begin{equation}\label{eq:v5-treebounds}
 a_0a(w)\le a(wi)\le\rho_q a(w),\qquad
 \sum_{i=1}^{b}a(wi)\le\rho_q a(w).
\end{equation}
\end{lemma}

\begin{proof}
Let $H$ bound the $\gamma$-H\"older seminorms of $\log h_i'$. The
chain rule, read from the innermost factor, gives
\[
 |\log h_w'(x)-\log h_w'(y)|
 \le H\sum_{j\ge0}r_+^{j\gamma}\le H/(1-r_+^\gamma).
\]
This proves the first assertion with $D=\exp\{H/(1-r_+^\gamma)\}$.
The upper multiplicative bound follows from the chain rule; for the
lower bound evaluate the inner derivative at a maximizer and bound
the outer derivative below by $r_u/D$.

Put $n=|w|$. For $k<n$, the points $T^kx,T^ky$ belong to the same
suffix cylinder. Their distance is at most $r_{w_{k:}}$. Immediately
after that suffix they belong to two different first-level cylinders,
so the mean value theorem also bounds this distance below by
$g r_{w_{k:}}/D$. At $k=n$ their distance is at least $g$. The
remaining squared distances are at most one. Multiplying by
$(1-q)q^k$ and summing proves \eqref{eq:v5-orbitseparation}, including
the term $q^n$ by means of its first summand $(1-q)q^ng^2$.
The same upper bound $V(w)$ holds for any two points of $K_w$.

Split the defining sum for $V(uv)$ at $|u|$ and apply
\eqref{eq:v5-distortion} to its first part to obtain
\eqref{eq:v5-concat}. Each suffix derivative $r_{v_{k:}}$ is
at least $r_v$, since the omitted outer maps are contractions.
Consequently $V(v)\ge r_v^2$. Equation~\eqref{eq:v5-concat}
therefore gives $V(uv)\le V(u)V(v)$. Since $V(u)\le1$ and
shift invariance gives
\[
 \mu[uv]\le\mu(S^{-|u|}[v])=\mu[v],
\]
we obtain the exact suffix domination $a(uv)\le a(v)$. No
independence of the symbols is used.

For a one-symbol child, $V(i)=q+(1-q)r_i^2\le\rho_q$;
submultiplicativity gives $V(wi)\le\rho_q V(w)$. Conversely,
split at $|w|$ and apply the lower distortion estimate to the
old summands. Since $V(i)\ge r_i^2\ge r_-^2$,
\[
 V(wi)\ge D^{-2}r_-^2\bigl(V(w)-q^{|w|}\bigr)
                  +r_-^2q^{|w|}
       \ge D^{-2}r_-^2V(w).
\]
The Gibbs bounds imply a positive lower bound $\beta$ on
$\mu[wi]/\mu[w]$, uniformly in $w,i$: divide the lower
bound for $[wi]$ by the upper bound for $[w]$ at the same
point, and use $\varphi\ge\inf\varphi$. Take
$a_0=\beta D^{-2}r_-^2/2$, decreased if necessary so that
$a_0<\rho_q$. Finally the children partition $[w]$; summing
$\mu[wi]V(wi)\le\rho_q\mu[wi]V(w)$ proves the final assertion.
\end{proof}

\begin{lemma}[Balanced greedy partitions]\label{lem:v5-greedy}
Let $L_N=|\mathcal P_N|$. Then $N/b\le L_N\le N$.
For $N\ge b$, let $\theta_N$ be the energy of the last leaf split
in constructing $\mathcal P_N$. Every $w\in\mathcal P_N$ satisfies
\begin{equation}\label{eq:v5-balanced}
 a_0\theta_N\le a(w)\le\theta_N,
 \qquad a_0L_N\theta_N\le G_N\le L_N\theta_N.
\end{equation}
For every fixed $A\ge1$ there is $c_A>0$ such that
\begin{equation}\label{eq:v5-dilation}
 c_AG_N\le G_{\lceil AN\rceil}\le G_N
 \quad(N\ge1).
\end{equation}
Moreover, for each $c>0$, the number of all finite words $v$ with
$a(v)\ge c\theta_N$ is at most $C_cN$.
\end{lemma}

\begin{proof}
The leaf counts are $1+j(b-1)$. A maximal energy selected for
splitting never increases, since each new child has at most $\rho_q$ times
its parent's energy. A leaf still present at the last split was
born from a parent of energy at least $\theta_N$, and every current
leaf has energy at most $\theta_N$. This proves
\eqref{eq:v5-balanced}; the root-only cases are kept separate.
The sum $G_N$ is nonincreasing by \eqref{eq:v5-treebounds}.

For completeness, the reverse comparison in \eqref{eq:v5-dilation}
does not require a preassigned power law. Choose an integer $j$ with
$b^{j+1}>b(A+1)$. If
$\theta_{\lceil AN\rceil}<a_0^{j+1}\theta_N$, every descendant
through depth $j$ of each old leaf would have been split:
all these nodes have energy at least $a_0^{j+1}\theta_N$.
The new partition would then have at least $b^{j+1}L_N$
leaves, which exceeds $\lceil AN\rceil$. Thus
$\theta_{\lceil AN\rceil}\ge a_0^{j+1}\theta_N$.
Combining this with \eqref{eq:v5-balanced} proves the reverse
comparison for $N\ge b$. Finitely many smaller budgets can be
absorbed into $c_A$; their energies are positive.

The finite full tree ending at $\mathcal P_N$ has at most
$bL_N/(b-1)$ vertices. Every other vertex descends from a leaf.
At depth $j$ below a leaf its energy is at most
$\rho_q^j\theta_N$. Only a bounded number of additional
\emph{generations} can therefore have energy at least
$c\theta_N$; each has at most $b^j$ descendants per leaf.
This bounds the total number of such words by $C_cN$.
\end{proof}

\begin{lemma}[Quantization of the orbit image]\label{lem:v5-quantization}
There are $0<c_2\le C_2<\infty$ such that
\[
 c_2G_M\le e_M^2(Z_q)\le C_2G_M\qquad(M\ge1).
\]
The lower bound allows arbitrary off-image centres in $\ell^2$.
\end{lemma}

\begin{proof}
Choosing one orbit point from each cylinder of $\mathcal P_M$
uses at most $M$ centres and has squared error at most $G_M$,
by the cylinder diameter estimate.

For the reverse bound take $\mathcal P_{4bM}$. It has at least
$4M$ leaves. For each leaf $w$ let
\[
 U_w=\{z\in\ell^2:\dist(z,Z_q(K_w))<\delta\sqrt{V(w)}\},
 \qquad 0<\delta<\sqrt{c_1}/(2\sqrt2).
\]
These open sets are pairwise disjoint. Indeed, two different leaves
have a common prefix $u$ strictly shorter than either leaf. Any
points in their two orbit cylinders have distance at least
$\sqrt{c_1V(u)}$, whereas $V(w),V(v)\le2V(u)$.
The sum of the two tube radii is smaller than this separation.
Thus $M$ centres meet at most $M$ of the tubes. Every unoccupied
cylinder contributes at least $\delta^2a(w)$ to the distortion.
By \eqref{eq:v5-balanced}, the sum over the unoccupied leaves is
at least a fixed multiple of $G_{4bM}$. Lemma~\ref{lem:v5-greedy}
compares this with $G_M$. The argument uses distances to entire
cylinders, so it does not require any centre to belong to the
orbit image. The finitely many root-only cases can again be
absorbed: the measure is nonatomic and the orbit's first
coordinate separates distinct states, so each fixed finite
quantization error is positive.
\end{proof}

\begin{lemma}[Exact suffix closure of the greedy tree]\label{lem:v5-suffix}
Let $\mathcal T_N$ be the set of all vertices, including the root,
in the greedy tree with leaf set $\mathcal P_N$, and write
$L_N=|\mathcal P_N|$. Then $\mathcal T_N$ is closed under deletion
of the first letter and
\begin{equation}\label{eq:v5-exact-states}
 |\mathcal T_N|=\frac{bL_N-1}{b-1}\le\frac{bN-1}{b-1}.
\end{equation}
In particular, every suffix of every leaf is already a vertex of
this same tree. No additional age states are required.
\end{lemma}

\begin{proof}
Full support and stationarity give, for every nonempty $u$,
\[
 \mu[uv]<\sum_{|z|=|u|}\mu[zv]=\mu[v].
\]
Together with $V(uv)\le V(v)$, this yields the strict inequality
$a(uv)<a(v)$. Suppose a vertex $w$ is selected for splitting.
Every proper suffix $v$ of $w$ was already an internal vertex at
that time. Otherwise some prefix $z$ of $v$ would be a leaf of
the current tree, and the decreasing child energies would give
$a(z)\ge a(v)>a(w)$, contradicting the maximality of $w$.
Thus the internal vertices are suffix closed.

If $w$ is a leaf, the parent of each nonempty proper suffix of $w$
is a suffix of the parent of $w$. This parent is internal by the
preceding argument (with the empty word interpreted as the root).
Its children are vertices of the tree, so the suffix itself is
also a vertex. The root and root-only tree cause no exception.
This proves closure for all vertices. Finally, a finite full
$b$-ary tree with $I$ internal vertices has $bI$ edges and $I+L_N$
vertices. Hence $bI=I+L_N-1$, giving \eqref{eq:v5-exact-states}.
\end{proof}

\begin{proof}[Proof of Theorem~\ref{thm:v5-intrinsic}]
The posterior-orbit converse, Theorem~\ref{thm:v4-orbit}, applied
to exact observations and geometric cycles, gives
$e_M^2(Z_q)\le R_M$, with a liminf bound for each observer.
Its randomization representation is Lemma~\ref{lem:v4-kernels}.
Equivalently, at each acquisition one freezes independent coding
randomness and the past; the $M$ possible post-update states
supply at most $M$ no-acquisition decoder strings. Quantization
bounds their weighted squared loss below, and the disjoint
last-acquisition events give the time-average bound. This does
not restrict the competing observers to symbolic codes.

For achievability, use $\mathcal T_N$ as the state set. At an acquisition $x$, write its unique leaf word.
On a nonacquisition update delete its first letter; the empty
word is absorbing until the next acquisition. Decode a nonempty
word $v$ by $\pi(v111\cdots)$ and the empty word by
$\pi(111\cdots)$. At age $k<|w|$ after acquisition in $K_w$,
the squared error is at most $r_{w_{k:}}^2$; after exhaustion it
is at most one. The average squared loss is at most $G_N$.
Every piece of information used by this rule is in its suffix
state. Neither reset age nor a precision index is supplied from
outside the register.

By Lemma~\ref{lem:v5-suffix}, this machine uses exactly
$(bL_N-1)/(b-1)$ designated states and at most $C_0N$ states, with
$C_0=b/(b-1)$. Take $N=\lfloor M/C_0\rfloor$ for sufficiently large $M$.
The fixed-dilation comparison gives $G_N\le CG_M$. For the
remaining finitely many budgets use the one-state constant
observer and enlarge $C$. Lemma~\ref{lem:v5-quantization} proves
the rest of \eqref{eq:v5-equivalence}.

Let $d_k$ be the expected age-$k$ loss of the constructed machine
under one acquisition. At time $t$ its expected loss is
\[
 \sum_{k<t}(1-q)q^kd_k+q^td_t
 \le\frac{1}{1-q}\sum_{k\ge0}(1-q)q^kd_k.
\]
The right side is at most $G_N/(1-q)$. The supremum objective
is bounded below by the average objective, proving its stated
order. When $q$ stays in a compact subinterval of $(0,1)$ and the
branch and Gibbs data are controlled, $\rho_q$ stays below one
and the separation and child-energy constants have common bounds.
All constants in the proof can then be chosen uniformly.
\end{proof}
